% ============================================================
%  Příprava na maturitu – Mechatronika
%  Kompiluj: pdflatex maturita_mechatronika.tex
%  Kódování: UTF-8
% ============================================================

\documentclass[a4paper,10pt]{article}

% --- Balíčky ---
\usepackage[T1]{fontenc}
\usepackage[utf8]{inputenc}
\usepackage[czech]{babel}
\usepackage[left=2.4cm, right=1.8cm, top=1.8cm, bottom=1.8cm]{geometry}
\usepackage{lmodern}
\usepackage{microtype}
\usepackage{parskip}
\usepackage{titlesec}
\usepackage{enumitem}
\usepackage{xcolor}
\usepackage{hyperref}
\usepackage{tabularx}
\usepackage{booktabs}
\usepackage{array}
\usepackage{tikz}
\usepackage{eso-pic}
\usepackage{mdframed}
\usepackage{amsmath}
\usepackage{amssymb}
\usepackage{pgfplots}
\pgfplotsset{compat=1.18}
\usetikzlibrary{arrows.meta, positioning, shapes.geometric, calc}

% --- Barvy ---
\definecolor{accent}{HTML}{03254E}
\definecolor{stripeblue}{HTML}{568EA3}
\definecolor{medgray}{HTML}{888888}
\definecolor{lightblue}{HTML}{EAF2F8}
\definecolor{lightyellow}{HTML}{FDFBE4}
\definecolor{lightgreen}{HTML}{EAF7EA}

% --- Modrý pruh vlevo ---
\AddToShipoutPictureBG{%
  \begin{tikzpicture}[remember picture, overlay]
    \fill[stripeblue]
      (current page.north west)
      rectangle
      ([xshift=10mm] current page.south west);
    \fill[accent!60]
      ([xshift=10mm] current page.north west)
      rectangle
      ([xshift=12mm] current page.south west);
  \end{tikzpicture}%
}

\hypersetup{colorlinks=true, urlcolor=accent, linkcolor=accent}

\titleformat{\section}
  {\large\bfseries\color{accent}}
  {\thesection.}
  {0.5em}
  {}
  [\color{accent}\titlerule]

\titlespacing{\section}{0pt}{16pt}{6pt}

\newmdenv[
  backgroundcolor=lightblue,
  linecolor=stripeblue,
  linewidth=1pt,
  roundcorner=4pt,
  innerleftmargin=8pt,
  innerrightmargin=8pt,
  innertopmargin=6pt,
  innerbottommargin=6pt,
  skipabove=4pt,
  skipbelow=4pt
]{pojmybox}

\newmdenv[
  backgroundcolor=lightyellow,
  linecolor=accent!40,
  linewidth=1pt,
  roundcorner=4pt,
  innerleftmargin=8pt,
  innerrightmargin=8pt,
  innertopmargin=6pt,
  innerbottommargin=6pt,
  skipabove=4pt,
  skipbelow=8pt
]{vysvetlenibox}

\newmdenv[
  backgroundcolor=lightgreen,
  linecolor=accent!30,
  linewidth=1pt,
  roundcorner=4pt,
  innerleftmargin=8pt,
  innerrightmargin=8pt,
  innertopmargin=6pt,
  innerbottommargin=6pt,
  skipabove=4pt,
  skipbelow=8pt
]{vzorcebox}

\pagestyle{plain}

% ============================================================
\begin{document}

% --- Záhlaví ---
\begin{minipage}[t]{0.75\textwidth}
    {\Huge\bfseries\color{accent} Příprava na maturitu}\\[4pt]
    {\large\color{stripeblue} Mechatronika}\\[2pt]
    \textcolor{medgray}{\small Suchánek Lukáš \quad SPŠ Prosek \quad 2026}
\end{minipage}
\hfill
\begin{minipage}[t]{0.22\textwidth}
    \raggedleft
    \begin{tikzpicture}
        \draw[color=stripeblue, rounded corners=4pt, line width=1pt]
            (0,0) rectangle (3,1.2)
            node[midway, align=center, color=accent, font=\small\bfseries]
            {Otázek: 22};
    \end{tikzpicture}
\end{minipage}

\vspace{6pt}
\noindent\color{accent}\rule{\linewidth}{1.5pt}
\vspace{2pt}

% ============================================================
\section{Průmyslové roboty a manipulátory}

\begin{pojmybox}
\textbf{Klíčové pojmy:}
DOF, kinematický řetězec, přímá/inverzní kinematika, SCARA,
delta robot, kolaborativní robot (cobot), TCP, trajektorie,
Denavit-Hartenberg, singularita, pracovní prostor
\end{pojmybox}

\begin{vysvetlenibox}

\textbf{Co je průmyslový robot}\\
Průmyslový robot je programovatelný manipulátor, který automaticky
provádí pohybové úlohy (např. svařování, montáž, balení).
Skládá se z mechanické části (rameno), pohonů (motory), senzorů
a řídicí jednotky.

Základní vlastnost robota je \textbf{stupňů volnosti (DOF)}.

\textbf{Stupně volnosti (DOF)}\\
DOF říká, kolik nezávislých pohybů robot má.
Typický průmyslový robot má \textbf{6 DOF}:

\begin{itemize}[noitemsep, leftmargin=1.4em]
    \item 3 DOF pro polohu (X, Y, Z)
    \item 3 DOF pro orientaci (naklopení a otočení nástroje)
\end{itemize}

Díky tomu dokáže robot umístit nástroj do libovolné polohy a směru,
podobně jako lidská ruka.

\textbf{Kinematika robota}\\
Kinematika řeší pohyb robota bez ohledu na síly.

\begin{vzorcebox}

\textbf{Přímá kinematika}\\
Známe úhly kloubů a chceme zjistit polohu nástroje (TCP):

\[
\mathbf{T} = A_1(\theta_1) \cdot A_2(\theta_2) \cdots A_n(\theta_n)
\]

Výstupem je poloha a orientace TCP (Tool Center Point).

\vspace{4pt}

\textbf{Inverzní kinematika}\\
Známe cílovou polohu TCP a hledáme úhly kloubů $\theta_i$.

Tento problém je složitější, protože:
\begin{itemize}[noitemsep]
    \item může existovat více řešení
    \item někdy neexistuje žádné řešení
    \item používají se numerické metody (iterace, Jacobian)
\end{itemize}

\end{vzorcebox}

\textbf{Co je TCP}\\
TCP (Tool Center Point) je bod na konci robota, který se sleduje
a programuje. Například špička svařovacího hořáku nebo uchopovač.

\textbf{Pracovní prostor robota}\\
Je to objem, kam robot dosáhne.
Závisí na:
\begin{itemize}[noitemsep]
    \item délce ramen
    \item konstrukci (SCARA, delta, artikulovaný)
    \item omezeních kloubů
\end{itemize}

\textbf{Typy průmyslových robotů}

\begin{center}
\begin{tabularx}{0.95\linewidth}{@{} l X l @{}}
\toprule
\textbf{Typ} & \textbf{Charakteristika} & \textbf{Použití} \\
\midrule

Artikulovaný (6-DOF) &
Nejvíce podobný lidské ruce, velmi flexibilní,
umožňuje složité pohyby ve 3D prostoru &
Svařování, lakování, manipulace \\

SCARA (4-DOF) &
Rychlý pohyb v rovině XY, velmi tuhý ve vertikále,
ideální pro montáž &
Montáž elektroniky, pick\&place \\

Delta robot (3-DOF) &
Paralelní konstrukce, extrémně rychlý a lehký,
malá hmotnost pohyblivých částí &
Balení, farmaceutický průmysl \\

Kartézský robot (3-DOF) &
Pohyby v osách X, Y, Z (lineární vedení),
jednoduchá konstrukce &
CNC stroje, 3D tisk \\

Cobot (kolaborativní robot) &
Navržen pro práci s člověkem,
má omezení síly a rychlosti pro bezpečnost &
Montáž, asistovaná výroba \\

\bottomrule
\end{tabularx}
\end{center}

\textbf{Trajektorie pohybu}\\
Trajektorie říká, jak se robot pohybuje mezi body:

\begin{itemize}[noitemsep, leftmargin=1.4em]
    \item \textbf{PTP (Point-to-Point)} -- robot jede nejrychlejší cestou,
    dráha mezi body není důležitá
    \item \textbf{CP (Continuous Path)} -- důležitá je přesná dráha
    (např. svařování nebo lakování)
    \item \textbf{Trapézový profil rychlosti} -- pohyb má 3 fáze:
    zrychlení, konstantní rychlost, zpomalení
\end{itemize}

\textbf{Singularita (důležitý pojem)}\\
Singularita nastává, když robot ztrácí jeden stupeň volnosti
a nemůže se pohybovat určitým směrem.
V praxi to může způsobit:
\begin{itemize}[noitemsep]
    \item náhlé zrychlení kloubů
    \item ztrátu kontroly trajektorie
\end{itemize}

\end{vysvetlenibox}

% ============================================================

% ============================================================
\end{document}
